% ===================================================================
%  TABELLA nr. 9 — La muntogna. Il guaud.
%  Traduction 2026 d'apres texto/io, controlee sur texto/fr, texto/ca
%  et texto/oc. Consigne : voir texto/rm/10-tabella-01.tex.
%
%  HUIT TABLEAUX DURANT, LE ROMANCHE A TRADUIT UN MONDE DE PLAINE ; SUR
%  CELUI-CI L'EMPRUNT CHANGE DE SENS.
%
%  Les mots dont cette page a besoin pour la haute montagne ne sont
%  pas des traductions de mots francais. Plusieurs sont ce dont les
%  mots francais sont faits :
%
%    « piz »        le sommet. Le mot n'a pas ete traduit sur les
%                   cartes d'Europe : Piz Bernina, Piz Palü, Piz
%                   Kesch se lisent tels quels dans toutes les
%                   langues ;
%    « lavina »     l'avalanche. Le francais « avalanche » vient
%                   d'une forme romane alpine du meme type, remontee
%                   des vallees vers la plaine ;
%    « muntanella » la marmotte. « Marmotte » sort de la meme
%                   fabrique alpine, le rat de montagne ;
%    « vadretg »    le glacier. Ni l'italien « ghiacciaio » ni le
%                   francais « glacier » ne lui ressemblent : ils
%                   derivent tous deux du mot « glace », le romanche
%                   a nomme la chose autrement ;
%    « crap »       le rocher, mot d'avant le latin, du fonds alpin ;
%    « chamutsch »  le chamois.
%
%  LE TABLEAU 5 DISAIT : LE CONTACT OFFRE, C'EST LE MANQUE QUI PREND.
%  LE TABLEAU 9 DIT L'AUTRE MOITIE : C'EST CELUI QUI A L'OBJET QUI
%  DONNE LE MOT. Le romanche a pris « chegls » et « schrinari » au
%  nord parce que le jeu et l'atelier venaient de la ; la plaine lui a
%  pris « lavina », « muntanella » et « piz » parce que la montagne
%  etait a lui. Les deux mouvements obeissent a la meme regle, et
%  cette regle n'a rien a voir avec le prestige d'une langue : elle
%  regarde qui tient la chose.
%
%  C'EST DONC ICI, ET ICI SEULEMENT DE TOUTE LA COLONNE, QUE LE
%  ROMANCHE NE TRADUIT PAS. Il rend son propre pays, et le livret
%  francais de 1926 employait deja, sans le savoir, trois de ses mots.
%
%  UNE INVERSION, LA PREMIERE DE LA COLONNE, AU BLOC t09-c1-16-1.
%  L'ido range guide (48), sac (47), touriste (46), rochers (45), et
%  la premiere redaction romanche mettait l'attaque des rochers avant
%  le sac : 48, 45, 47, 46. renvoji.py l'a signalee. La phrase a ete
%  refaite en portant les deux complements de maniere devant le verbe
%  — « cun in satg sin il dies e cun l'alpenstock en maun, sco in ver
%  turist, assagl las crappas » —, ce qui rend l'ordre de l'ido sans
%  rien forcer. La colonne reste a zero divergence.
% ===================================================================

%%K t09-apar-1 apar
\VUsaut{4.0mm}
\VUtitre{15.0pt}{60}{TABELLA nr. 9}
\VUsaut{3.0mm}

%%K t09-tit-1 sub
\VUcentre{12.0pt}{0}{\textit{Emprima scena.}}
\VUsaut{3.0mm}

%%K t09-tit-2 sub
\VUpk{11.4pt}{40}{La muntogna. --- La citad da bogns.}
\VUsaut{3.0mm}

%%K t09-c1-01-1 p
1. --- Dapi otg dis sun jau a Pratz. Jau na poss betg exprimer tut l'admiraziun
ch'jau hai sentì cur ch'jau sun stà avant la maravigliusa \VUgras{chaglia da muntognas}
\textsuperscript{(1)} da las Pireneas, da la quala la \VUgras{cresta}
\textsuperscript{(2)} para sin il tschiel blau sco in feston gigantic.

%%K t09-c1-02-1 p
2. --- Jau na m'imaginava betg ch'las \VUgras{cimas} \textsuperscript{(3)} pireneicas
cuntanschian ina \VUgras{autezza} uschè gronda, e jau sun restà stupefatg avant ils
bellissims \VUgras{vadretgs} \textsuperscript{(5)} ed ils \VUgras{pizs}
\textsuperscript{(4)} da naiv, sin ils quals rullan \VUgras{lavinas} uschè terrorisantas.

%%K t09-tit-3 sub
\VUsaut{3.0mm}
\VUpk{10.2pt}{40}{Cuntrada da muntogna.}
\VUsaut{3.0mm}

%%K t09-c1-03-1 p
3. --- Gia il di suenter mia arrivada sun jau ì al vegl \VUgras{vulcan}
\textsuperscript{(6)} stizzà ch'è sper Pratz. Sin ils urs arids e nivs dal
\VUgras{crater} vesa ins anc bain las trattas dal passadi da la \VUgras{lava} ch'è
currida, avant pliras tschientaners, sin il \VUgras{flanc da la muntogna}
\textsuperscript{(7)}. Jau hai cuntanschì da chattar, sin in dals \VUgras{pendents},
intgins toctgs da quella lava.

%%K t09-c1-04-1 p
4. --- Pratz sa chatta sin la frontiera, e la \VUgras{fortezza} \textsuperscript{(8)} che
defenda la \VUgras{stretga} \textsuperscript{(27)} che separa la Frantscha da la Spagna fa
propi s'algordar da quels vegls chastels da la riva dal Rain che paran spiar ina
\VUgras{preda} da sisum lur grep.

%%K t09-c1-05-1 p
5. --- In'enorma \VUgras{evla} \textsuperscript{(9)} ch'ha ses gnieu betg lunsch dal
\VUgras{fort} \textsuperscript{(8)} tira savens mia attenziun.

Quel \VUgras{utschè da rapina}, dal qual il \VUgras{pichel} è ferm, porta savens a ses
pitschens in agnè u in chavrin ch'el tegna cun sias \VUgras{unglas}. Uschè gronda è la
grondezza da sias \VUgras{alas} ch'el po sgular ditg senza sbatter cun sia chargia
pesanta.

%%K t09-tit-4 sub
\VUsaut{3.0mm}
\VUpk{10.2pt}{40}{L'excursiunist.}
\VUsaut{3.0mm}

%%K t09-c1-06-1 p
6. --- Là è la pli agreabla spassegiada l'excursiun a la cascada \textsuperscript{(10)}
da «Cau». La \VUgras{crudada d'aua} croda en vortischs e sparescha lura sco in bel
\VUgras{auatg} en in \VUgras{guaud da pignels} \textsuperscript{(11)} ch'è a vest da la
citad. Nus essan muntads tras quest guaud fin a l'\VUgras{observatori meteorologic}
\textsuperscript{(12)}, e da sisum il \VUgras{belveder} avain nus survesì tut il
\VUgras{panorama} da la regiun. La \VUgras{cuntrada} è maravigliusa, e la
\VUgras{natira} para prodigar al pajais tut ils tresors ch'ella ha.

%%K t09-c1-07-1 p
7. --- Per ir giu avain nus duvrà il \VUgras{funicular} \textsuperscript{(13)} e nus
essan restads en ina pitschna \VUgras{staziun} sper las \VUgras{ruinas}
\textsuperscript{(14)} d'in vegl \VUgras{chastè feudal}, abità ina giada dals signurs da
Pratz.

%%K t09-c1-08-1 p
8. --- Paupra chastè! Tge patratga el forsa vesend giu la coketta \VUgras{citad da
bogns} \textsuperscript{(15)}, ch'è ussa ina \VUgras{staziun termala} a la moda?

%%K t09-c1-08-2 p
Il \VUgras{clima} è uschè agreabel, l'\VUgras{aua minerala} uschè efficazia, ch'la
\VUgras{cura} duvrada en il \VUgras{sanatori} \textsuperscript{(17)} è in ver plaschair.

%%K t09-tit-5 sub
\VUsaut{3.0mm}
\VUpk{10.2pt}{40}{Agreabla citad termala.}
\VUsaut{3.0mm}

%%K t09-c1-09-1 p
9. --- Dal rest ha ins fatg tut per render agreabel il suggiurn. In grond
\VUgras{viaduct} \textsuperscript{(16)} è vegnì construì per pussibilitar ina viafierra;
il \VUgras{parc} \textsuperscript{(18)} dal \VUgras{stabiliment termal}, en il qual ins
\VUgras{fa concerts}, è arrangià en ina moda charmanta. Pliras graziusas
«\VUgras{villas}» \textsuperscript{(19)} èn vegnidas construidas enturn il casino
\textsuperscript{(21)}, ed ina \VUgras{chaussée} \textsuperscript{(22)} fitg bain
mantegnida facilitescha ferm las communicaziuns.

%%K t09-c1-10-1 p
10. --- Ina simpla \VUgras{scarpa} \textsuperscript{(23)} separa la \VUgras{via}
\textsuperscript{(20)} da la \VUgras{val} \textsuperscript{(25)} propriamain ditga, en il
fund da la quala sa chatta in grazius pitschen \VUgras{lai} \textsuperscript{(26)}
ch'alimentescha
in \VUgras{auatg} \textsuperscript{(24)} limpid.

%%K t09-c1-11-1 p
11. --- Da l'autra vart da la val vegn explotada ina \VUgras{minera da fier}
\textsuperscript{(28)}. Pliras giadas sun jau ì a vesair ils \VUgras{minaders} che van
giu en il \VUgras{poz}, e jau sun schizunt entrà en la \VUgras{galaria} en la quala els
passentan lur di, extrahend il \VUgras{mineral} che cuntegna il \VUgras{metal}.

%%K t09-c1-12-1 p
12. --- Ina impurtanta \VUgras{fusina} è vegnida construida sper la minera per
utilisar immediatamain las ritgezzas ch'ins extrahia da ella. Il \VUgras{fuorn aut}
\textsuperscript{(29)}, stgaudà cun il \VUgras{carvun da terra} ch'i dat qua a bler,
pussibilitescha d'obtegnair svelt e bunmartgà \VUgras{fier fusì}.

%%K t09-c1-13-1 p
13. --- Betg lunsch da la minera vegn chavada ina \VUgras{chava da crappa}
\textsuperscript{(30)}. Ni \VUgras{granit}, ni \VUgras{porfir}, ni gnanc \VUgras{crappa da
sablun} tagliescha il vegl \VUgras{chavader} cun ses \VUgras{pichet}, mabain simpla
\VUgras{crappa da taglia} per la construcziun da chasas.

%%K t09-c1-14-1 p
14. --- In dals pli bels ornaments da quel pajais èn ils \VUgras{pignels}
\textsuperscript{(31)}. Dapertut en il fund d'ina \VUgras{ruvina} \textsuperscript{(32)},
sin la riva d'in \VUgras{torrent} \textsuperscript{(33)}, d'in \VUgras{abiss}
\textsuperscript{(34)} u d'in precipizi, dattan lur sutglas guaglias ina nianza verda.

%%K t09-tit-6 sub
\VUsaut{3.0mm}
\VUpk{10.2pt}{40}{Ir a pe.}
\VUsaut{3.0mm}

%%K t09-c1-15-1 p
15. --- Jau profitesch da mias vacanzas per \VUgras{spassegiar ditg}. Las
\VUgras{vias} \textsuperscript{(35)} che sinueschan sin la muntogna pussibiliteschan da far,
senza dar quità a la distanza, pliras \VUgras{kilometers} e savens plirs
\VUgras{miriameters}.

%%K t09-c1-15-2 p
\VUgras{Termins} \textsuperscript{(36)} e \VUgras{pals indicaturs} \textsuperscript{(37)}
marcheschan las vias, sin las qualas ins scuntra savens in giuven \VUgras{muntagnard}
\textsuperscript{(38)} cun ina \VUgras{bisatga} \textsuperscript{(40)} sper il flanc, che
porta ina \VUgras{muntanella} \textsuperscript{(39)}, u in \VUgras{zigan}
\textsuperscript{(41)}, da la quala la \VUgras{giacca da pelitscha} \textsuperscript{(42)}
fa in contrast curius cun la veglia \VUgras{capuscha} \textsuperscript{(43)} sfessa. El
maina cun ina chadaina in \VUgras{urs} \textsuperscript{(44)} dumagnà.

%%K t09-tit-7 sub
\VUpk{10.2pt}{40}{Muntada sin las muntognas.}
\VUsaut{3.0mm}

%%K t09-c1-16-1 p
16. --- Quasi mintga di vom jau cun in \VUgras{guid} \textsuperscript{(48)} a far ina
\VUgras{muntada}, e jau sun fitg fiers cur ch'jau, cun in \VUgras{satg}
\textsuperscript{(47)} sin il dies e cun l'«\VUgras{alpenstock}» en maun, sco in ver
\VUgras{turist} \textsuperscript{(46)}, assagl las \VUgras{crappas}
\textsuperscript{(45)}.

%%K t09-c1-17-1 p
17. --- Da sisum ina muntogna paran ils abitants da la plaunira betg pli gronds che
furmiclas; ina \VUgras{scossa da nursas} \textsuperscript{(50)} che pascenta en ina
\VUgras{pastira} \textsuperscript{(49)} para in giugarel per uffants. In \VUgras{pign}
\textsuperscript{(51)} u era ina \VUgras{chasetta da lain} \textsuperscript{(52)} para apaina
ina maglia sin il verd.

%%K t09-c1-18-1 p
18. --- Durant quellas excursiuns scuntrain nus savens sin il \VUgras{truoi}
\textsuperscript{(53)} in \VUgras{chatschader da chamutschs} \textsuperscript{(54)} che porta
sin la spatla la bes-cha ch'el ha gist mazzà.

%%K t09-c1-19-1 p
19. --- Jau hai mess en mes erbari in rom fitg remarcabel da \VUgras{filesch}
\textsuperscript{(55)} ed in bel rometg rosa da \VUgras{bruch} \textsuperscript{(57)},
ch'jau hai racoltà a l'entrada d'ina pitschna \VUgras{grotta} \textsuperscript{(56)}
durant mia ultima spassegiada.

%%K t09-tit-8 sub
\VUsaut{3.0mm}
\VUcentre{12.0pt}{0}{\textit{Segunda scena.}}
\VUsaut{3.0mm}

%%K t09-tit-9 sub
\VUpk{11.4pt}{40}{Il guaud. --- La chatscha.}
\VUsaut{3.0mm}

%%K t09-c2-01-1 p
1. --- Ier hai jau passentà in di delizius en il guaud. La diligenza che regia
tranter ils animals e las \VUgras{inserts} \textsuperscript{(5)} che l'abitan m'ha fitg
interessà.

%%K t09-c2-01-2 p
L'\VUgras{arogna} \textsuperscript{(1)} che tessa sia \VUgras{rait} \textsuperscript{(2)}
tranter dus roms per tschiffar la \VUgras{mustga} \textsuperscript{(3)} che passa, la
\VUgras{chaglia} \textsuperscript{(4)} che porta cun grondas fadias sia chascha, il
scarvatsch che explorescha sia preda, la diligenta furmicla fitg zelusa sper la
\VUgras{furmiclera} \textsuperscript{(6)}, tut quels pitschens èn ids, vegnids, han lavurà
cun ina diligenza extraordinaria.

%%K t09-c2-02-1 p
2. --- Ina \VUgras{serp} \textsuperscript{(7)}, ch'jau hai a l'emprim vis tegnì per ina
\VUgras{vipra}, ma che n'era nagut auter ch'in simpel \VUgras{culiver}, sa zuppava sut
in chip d'arber e sportgava da temp en temp ses sutgl chau, che pareva adina pront
per morder. Avant mai vaseva plaun ina gronda \VUgras{luscha} \textsuperscript{(8)}
suenter avair devorà ina da las savurusas \VUgras{fravgias} \textsuperscript{(11)} che
creschan qua a bler.

%%K t09-c2-03-1 p
3. --- Il terren era tapezzà cun \VUgras{flurs da guaud}: \VUgras{primlas}
\textsuperscript{(9)}, \VUgras{pervincas} \textsuperscript{(10)} e bleras autras plantas
ch'jau n'enconuscheva betg fluriva tranter \VUgras{tuffs d'erva} \textsuperscript{(12)}.

%%K t09-c2-04-1 p
4. --- En quella regiun è la \VUgras{trifla} quasi uschè communa sco il
\VUgras{bulieu} \textsuperscript{(14)}, ed in \VUgras{purtgrin} \textsuperscript{(13)} che
tutga ad in selvicultur explorava gurmandamain sch'el vegniss a chattar intginas
d'ellas.

%%K t09-tit-10 sub
\VUsaut{3.0mm}
\VUpk{10.2pt}{40}{Chatschar da furtim.}
\VUsaut{3.0mm}

%%K t09-c2-05-1 p
5. --- Jau hai assistì a la \VUgras{fin} d'ina scena che m'ha fitg divertì. Cun agid
da l'udur da ses \VUgras{chaun bass} \textsuperscript{(15)}, ha il \VUgras{guardgiaboss}
\textsuperscript{(16)} gist surprendì in \VUgras{chatschader da furtim}
\textsuperscript{(22)}, en il mument che quel sa preparava a purtar davent la
\VUgras{selvaschina} \textsuperscript{(17)} ch'el aveva mazzà. Preferind da bandunar sia
preda enstagl da sa laschar arrestar, era il chatschader da furtim fugì, ed ina
\VUgras{levra} \textsuperscript{(18)}, ina \VUgras{cerva} \textsuperscript{(19)}, in
\VUgras{chavriel} \textsuperscript{(20)} ed in \VUgras{portg selvadi} \textsuperscript{(21)}
giaschevan sin l'erva, allegrond il guardgiaboss.

%%K t09-c2-06-1 p
6. --- La diversitad dals arbers ch'il guaud cuntegna è surprendenta. Ils
\VUgras{faus} \textsuperscript{(23)} e las \VUgras{betlas} \textsuperscript{(26)}, ils
\VUgras{pigns}, che dattan la \VUgras{rasa} \textsuperscript{(27)}, ils \VUgras{ruvers}
\textsuperscript{(29)} ed era blers auters mischeidan lur roms.

%%K t09-c2-06-2 p
L'\VUgras{eider} \textsuperscript{(24)}, che sa taca al chip dals auts arbers, dat al
\VUgras{guaud} \textsuperscript{(25)} ina verdezza splendurusa, che sa maridescha
agreablamain cun la nianza pli clera e pallid verda dal \VUgras{mustgel}
\textsuperscript{(31)}, en il qual il \VUgras{pitgalain} \textsuperscript{(30)} tschertga
inserts.

%%K t09-tit-11 sub
\VUsaut{3.0mm}
\VUpk{10.2pt}{40}{Cuntrada da guaud.}
\VUsaut{3.0mm}

%%K t09-c2-07-1 p
7. --- Ils vegls ruvers m'han incantà. Lur grossas \VUgras{ragischs}
\textsuperscript{(32)} tortas, mez ordadora dal terren, als dattan in splendid aspect
da fermezza e vigur, e jau ma dumond co ch'il \VUgras{possessur} \textsuperscript{(35)}
po sa decider da laschar als taglier e da als dar a la \VUgras{resgia}
\textsuperscript{(34)} dal \VUgras{bostgader} \textsuperscript{(33)}. Ils pauprs
\VUgras{chips} \textsuperscript{(36)}, spogliads da lur \VUgras{scorsa}
\textsuperscript{(37)} u fendids da l'\VUgras{ascha} \textsuperscript{(38)} dal
\VUgras{lavurant a di} \textsuperscript{(39)}, ma fan displaschair.

%%K t09-c2-08-1 p
8. --- Dasperas ha in vegl \VUgras{carbunier} \textsuperscript{(41)} preparà cun sia
\VUgras{pala} in \VUgras{mantun} \textsuperscript{(42)} da carvun, ed ina veglia purtava
in \VUgras{fasch} \textsuperscript{(40)} da \VUgras{lain mort} fatg da rometgs dals arbers
tagliads.

%%K t09-tit-12 sub
\VUsaut{3.0mm}
\VUpk{10.2pt}{40}{Chatscha.}
\VUsaut{3.0mm}

%%K t09-c2-09-1 p
9. --- Jau vuleva ir a pausar intgins mumaints en la \VUgras{chasa dal guardgiaboss}
\textsuperscript{(46)}, ma, cur ch'jau sun arrivà sper il \VUgras{laiet}
\textsuperscript{(43)}, sin il qual nuda in bel \VUgras{cign} \textsuperscript{(45)},
hai jau vis ch'in recent \VUgras{auadiv} \textsuperscript{(44)} aveva midà la via en in ver
\VUgras{palì}. Jau hai stuì far in gir e, passond sper il \VUgras{ceder}
\textsuperscript{(49)}, ils \VUgras{papels} \textsuperscript{(48)} e l'\VUgras{ulm}
\textsuperscript{(47)} ch'èn a l'ur dal guaud, hai jau vis sortir d'in \VUgras{bostg}
\textsuperscript{(54)} in bellissim \VUgras{tschierv} \textsuperscript{(50)}, persequità
d'ina obstinada \VUgras{muta da chauns} \textsuperscript{(51)}, ch'era era excitada dal
\VUgras{corn} \textsuperscript{(53)} dals \VUgras{chatschaders a chaval}
\textsuperscript{(52)}. La paupra bes-cha era dal tut extenuada. Ella è crudada
morind a intgins pass da mai.

%%K t09-c2-10-1 p
10. --- Il figl dal guardgiaboss, ch'il canaster da la \VUgras{chatscha a curriment}
aveva attrat, è sortì da la chasa e nus essan turnads en dus en il guaud. El aveva
vis ina \VUgras{gaza} \textsuperscript{(56)} sin in rom d'in vegl ruver. El pensava che
ses gnieu saja probablamain zuppà en il \VUgras{fegliom} \textsuperscript{(58)} ed el
vuleva prender il gnieu.

%%K t09-c2-10-2 p
Agil sco in \VUgras{scuriatel} \textsuperscript{(61)}, è el rampignà da \VUgras{rom}
\textsuperscript{(55)} a rom, ma el n'ha chattà nagut ed ha mo cuntanschì da far sgular
in vegl \VUgras{cornagl} \textsuperscript{(60)} ch'era sesì sin ina \VUgras{rama}
\textsuperscript{(57)}.

\VUsaut{4.0mm}
\VUfilet{22mm}
